%%
%% Hypothesentest
%% 2023 - 05 - 24 φ
%%

\subsection{Würfelexperiment}

%%
%Bearbeiten Sie mit der Klasse das Würfelexperiment im OLAT:

%\weblink{Würfelexperiment}{https://olat.bms-w.ch/auth/RepositoryEntry/6029794/CourseNode/107682584460080}



\subsubsection{\epsdice{3} werfen}
\TRAINER{Die Würfel werden vorab so sortiert, dass die Hälfte der
Klasse einen gezinkten Würfel mit doppelter \epsdice{3} erhalten wird.}

Werfen Sie den Würfel 29 Mal und notieren Sie das Ergebnis wie folgt:
\begin{itemize}
\item Wie oft kommt die \epsdice{3},
\item wie oft mehr als die \epsdice{2}?
\end{itemize}

\begin{bbwFillInTabular}{|c|p{50mm}|}\hline
$\epsdice{3}$ & \\\hline
$\epsdice{3},\epsdice{4},\epsdice{5}$ oder $\epsdice{6}$ & \\\hline
\epsdice{1} oder \epsdice{2} & \\\hline
\end{bbwFillInTabular}

\subsubsection{Ergebnis}

Wie oft ist die \epsdice{3} erschienen? ....................
\vspace{15mm}

Wie oft kam mehr als die \epsdice{2}? .......................

\newpage

\subsubsection{Unsere Resultate}

Füllen wir in der Klasse gemeinsam die folgende Tabelle aus:

\begin{bbwFillInTabular}{p{5cm}|p{5cm}|p{5cm}}
    & Würfel ist «fair» ($H_0$) & Würfel ist gezinkt ($H_1$)\\\hline
Würfel wird als «fair» gewertet (weniger als acht \epsdice{3}-er).&
& \TRAINER{Fehler 2. Art}\\\hline
Würfel wird als gezinkt gewertet (mehr als
sieben \epsdice{3}-er).& \TRAINER{Fehler erster Art}& \\\hline
\end{bbwFillInTabular}


\subsubsection{Nullhypothese / Alternativhypothese}
Angenommen, der Würfel ist «fair», so erwarten wir, dass jede Seite
etwa gleich häufig auftritt.

\begin{center}\textbf{Nullhypothese} oder {$H_0$}.\end{center}

Die \textbf{Nullhypothese} $H_0$ sagt, dass einzig der Zufall im Spiel war,
sonst nichts.

Wir \textbf{vermuten} jedoch, dass hier etwas nicht mit rechten Dingen
zugeht und formulieren die Alternativhypothese:

Die \textbf{Alternativhypothese} lautet hier \zB «Es kommt häufiger
die \epsdice{3} als wenn der bloße Zufall $\left(\frac{1}{6}\right)$
im Spiel ist».

\subsubsection{Rechsseitig / Linksseitig}
Es handelt sich um einen \textbf{rechtsseitigen} Test, da wir
vermuten, dass das Resultat häufiger auftritt als erwartet.

\subsubsection{Testgröße}
Die \textbf{Testgröße} wird typischerweise mit $X$ bezeichnet und soll
hier die Anzahl der geworfenen \epsdice{3}-ern angeben.

\subsubsection{Erwartungswert}
Der \textbf{Erwartungswert} $\mathbb{E}(X)$ der Nullhypothese bei
29-maligem Würfeln liegt bei $\frac{29}{6} \approx 4.8$. Einige haben
mehr als 4 mal die \epsdice{3} geworfen. Bis wo hin glauben wir noch
an einen «fairen» Würfel? Hier kommt Bernoulli ins Spiel:


\bbwCenterGraphic{12cm}{geso/stoch/img/wuerfeltabelle.png}

\subsubsection{Signifikanzniveau}
Mit 89.08 \% Wahrscheinlichkeit werfe ich weniger als \textbf{sieben}
mal eine \epsdice{3} mit einem normalen Würfel. Wer also mehr als
\textbf{sechs} mal eine \epsdice{3} geworfen hat, ist schon eher
\textbf{verdächtig}. Hier kommt der Begriff \textbf{Signifikanzniveau} ins
Spiel. Wir wollen erst «anklagen», wenn wir zu 95 \% sicher sind, dass
der Würfel gezinkt war. Unser Signifikanzniveau legen wir mehr oder
weniger willkürlich  bei $\alpha = 0.05 = 5 \%$ fest. Mit $95.53 \%$
Wahrscheinlichkeit werfe ich maximal \textbf{sieben} \epsdice{3}-er.

Somit gilt bei einem Signifikanzniveau von 5 \%
folgendes:

Wer mehr als \textbf{sieben} mal eine \epsdice{3} wirft, fällt bereits
unter die 5 \%, die wir noch «glauben».

\subsubsection{Fehler erster Art ($\alpha$-Fehler)}

Es kann nun durchaus passieren, dass ein «fairer» Würfel als «gezinkt»
kategorisiert wird. Dieses fälschliche Zurückweisen der Nullhypothese
wird als \textbf{Fehler erster Art} oder als \textbf{$\alpha$-Fehler}
bezeichnet.

Der \textbf{Fehler 2. Art} (= \textbf{$\beta$-Fehler}) wäre dann die
fälschliche Beibehaltung der Nullhypothese; also dass ein «gezikter»
Würfel als «fair» taxiert wird.

\subsubsection{Entscheidungsregel}
Gesucht ist die Anzahl \epsdice{3}-er, wo der «normale» Bereich
«kippt». Wir gehen von einem Signifikanzniveau von 95 \% aus. Das
heißt: Bis zu welcher Anzahl $k$ ist die Anzahl geworfener
\epsdice{3}-er noch mit 95 \% Wahrscheinlichkeit «normal»?

Mit 95 \% Wahrscheinlichkeit kommen also 0 bis k \epsdice{3}-er?

Unsere Verteilung zeigt, dass bei maximal sechs \epsdice{3}-ern immer
noch etwa 11 \% Irrtumswahrscheinlichkeit vorhanden ist. Anders
ausgedrückt: Wenn wir bei «wir glauben maximal 6 \epsdice{3}er»
bleiben, so wird aber mit ca 11 \% wahrscheinlichkeit mehr als 6 mal
die \epsdice{3} geworfen. Wir müssen also noch höher gehen.
Wenn wir die Entscheidungsregel bei $k=7$ ansetzen, so sagen wir: Es
ist noch gut wahrscheinlich maximal $k=7$ mal eine \epsdice{3} zu
erzielen. Mehr als sieben Mal hingegen unterschreitet unser 5 \%
Signifikanzniveau. Dies bringt und zum nächsten Begriff:

\subsubsection{Annahmebereich / Ablehnugngsbereich}

Der Annahmebereich der Nullhypothese liegt bei uns mit einem 5 \%
Signifikanzniveau also bei 0 bis 7. Der Ablehnungsbereich unserer
Nullhypothese liegt somit bei 8 bis 29 mal. 


\newpage






